\subsection{Code-as-Policy: Robots that Write their Own Skills}
\label{sec:code}

The organising contribution of this survey is to arrange code-as-policy methods not by
application or robot embodiment, but by \emph{how much of their own competence a system
produces at run time}. Concretely, we ask five questions in sequence: does the agent (i)
write control code at all, (ii) repair that code from execution feedback within a task,
(iii) remember validated code across tasks, (iv) search a population of programs rather
than following a single trajectory of repairs, and (v) close all of these into one
open-ended loop? Each ``yes'' is a rung on the ladder of Figure~\ref{fig:taxonomy_31}, and
the population thins sharply toward the top: of the thirteen zero-shot systems, only a
handful ever reach the combined feedback--memory--search regime. The novelty of the
survey is naming this progression and showing that the top rung is nearly empty.

\input{figures/fig_taxonomy_31}
% Figure — the self-improvement ladder (representative rung: §3.1.5 ASPIRE loop).
\begin{figure}[t]
\centering
\begin{tikzpicture}[
  node distance=6mm and 9mm,
  box/.style={draw=hidden-draw, rounded corners, fill=secAl, align=center,
              font=\footnotesize, minimum height=8mm, text width=15mm, inner sep=3pt},
  core/.style={box, fill=secA},
  io/.style={draw=none, font=\footnotesize\itshape},
  ar/.style={-{Stealth[length=2mm]}, darkgray, line width=0.7pt},
]
  \node[io] (task) {task};
  \node[box, fill=secA, right=of task] (actor) {actor\\ agent};
  \node[box, fill=secB, right=of actor] (eng) {exec.\\ engine (F)};
  \node[box, fill=secC, right=of eng] (mem) {skill\\ memory (M)};
  \node[box, fill=secD, right=of mem] (srch) {evol.\\ search (S)};
  \node[io, right=of srch] (out) {validated skill};

  \draw[ar] (task) -- (actor);
  \draw[ar] (actor) -- (eng);
  \draw[ar] (eng) -- (mem);
  \draw[ar] (mem) -- (srch);
  \draw[ar] (srch) -- (out);
  % feedback loops
  \draw[ar] (eng.south) to[out=-90,in=-90] node[below,font=\scriptsize]{traces} (actor.south);
  \draw[ar] (mem.north) to[out=90,in=90] node[above,font=\scriptsize]{skills $\rightarrow$ future tasks} (actor.north);
\end{tikzpicture}
\caption{The full self-improving loop (\S3.1.5): feedback (F) $+$ memory (M) $+$ search (S)
combine into one open-ended loop---the ASPIRE cell. VLAs (\S3.2) have no such loop.}
\label{fig:architectures}
\end{figure}


\subsubsection{Zero-shot synthesis}
\label{ssec:zeroshot}
The foundational move is to treat the language model as a \emph{program synthesiser} over
a fixed library of perception and control primitives. \emph{Code-as-Policies}
\citep{codeaspolicies} showed that an LLM prompted with an API and a natural-language
instruction can emit executable Python that composes those primitives, recursively
defining undefined functions and parameterising control with arithmetic and feedback
logic. \emph{ProgPrompt} \citep{progprompt} situates the generated program in the scene
by exposing available actions and objects as importable symbols, while \emph{VoxPoser}
\citep{voxposer} has the model write code that \emph{composes 3D value maps} for a motion
planner rather than calling fixed skills, loosening the dependence on a hand-authored API.
Subsequent work broadened the input modality and task horizon---\emph{RoboCodeX}
\citep{robocodex} conditions on multimodal observations, \emph{Instruct2Act}
\citep{instruct2act} maps instructions to perception--action code, and video- and
demonstration-conditioned variants such as \emph{RoboPro} \citep{robopro},
\emph{Demo2Code} \citep{demo2code} and \emph{Text2Motion} \citep{text2motion} recover
programs from richer context. \emph{Statler} \citep{statler} maintains an explicit world
state across steps, and \emph{TidyBot} \citep{tidybot} and \emph{ChatGPT for Robotics}
\citep{chatgptrobotics} demonstrate personalisation and prompt design for real hardware.
What unites this rung, and limits it, is that the program is written \emph{once}: there is
no channel from execution back to the synthesiser, so a plan that fails at run time simply
fails.

\subsubsection{Closed-loop self-repair}
\label{ssec:selfrepair}
The second rung adds a within-task feedback loop. \emph{Inner Monologue}
\citep{innermonologue} feeds success detectors, scene descriptions and human feedback back
into the model as language, letting it replan when a step fails. \emph{DoReMi}
\citep{doremi} detects misalignments between plan and execution and triggers recovery, and
\emph{REFLECT} \citep{reflect} builds a hierarchical, multimodal summary of past
interaction that an LLM queries to explain a failure and propose a correction. More recent
systems push the monitor into the perception loop---\emph{Code-as-Monitor}
\citep{codeasmonitor} compiles spatio-temporal constraints into vision-language checks, and
\emph{AHA} \citep{aha} trains a model to reason explicitly about failure modes. The novelty
of this rung is the closed loop, but it is a \emph{short} loop: corrections are discarded
at task boundaries, so nothing accumulates.

\subsubsection{Skill-library accumulation}
\label{ssec:skilllib}
The third rung makes the loop persistent by writing validated code into a growing library
that later tasks retrieve and compose. Because this ``memory'' axis is a full technique
family in its own right, we catalogue its systems once, in \S\ref{sec:skill}
(\S3.4.3)---\emph{Voyager} \citep{voyager} is the canonical example, curating an
ever-growing library of verified skill programs with when-to-apply guards, while related
methods distil language corrections into retrievable knowledge for future tasks~\citep{droc}. The key
distinction from \S3.1.2 is temporal scope: competence compounds \emph{across} tasks rather
than being rebuilt each time.

\subsubsection{Evolutionary program search}
\label{ssec:search}
The fourth rung replaces a single chain of repairs with population-based search over
programs. Rather than iteratively debugging one candidate, the agent maintains several,
scores them by execution, and mutates the best. \emph{CaP-X} \citep{capx} provides an
interactive gym and benchmark for exactly this style of coding agent, and \emph{RoboEvolve}
\citep{roboevolve} couples a planner and a learned simulator into a co-evolutionary loop
that mines near-miss failures to stabilise search. The novelty is the shift from
\emph{repair} to \emph{search}: exploration is explicit and parallel, which trades compute
for a better chance of escaping local failures.

\subsubsection{The full self-improving loop}
\label{ssec:fullloop}
The top rung combines all three mechanisms---execution feedback (F), a persistent skill
memory (M), and evolutionary search (S)---into one open-ended loop
(Figure~\ref{fig:architectures}). \emph{ASPIRE} \citep{aspire} instantiates this with a
robot execution engine that emits per-primitive multimodal traces, a skill library of
validated code with guards, and a population search over candidate programs, so that
experience compounds and transfers across tasks and embodiments. Its real-world sibling
\emph{ENPIRE} \citep{enpire} closes the same loop on physical hardware, using automatic
reset, parallel rollouts, and log-driven algorithm revision to drive a coding agent to
near-perfect success on dexterous tasks. This cell---feedback \emph{and} memory \emph{and}
search---is what no prior survey isolates, and it is nearly empty: it is the flag this
survey plants.
